\loai{Từ thông cực đại và suất điện động cực đại}
\begin{vidu} %[3P3-1.1B][Loai2] 
	Một khung dây dẫn phẳng dẹt hình chữ nhật có 500 vòng dây, diện tích mỗi vòng là $220\ cm^2$. Khung quay đều với tốc độ 50 vòng/giây quanh một trục đối xứng nằm trong mặt phẳng của khung dây, trong một từ trường đều có vector cảm ứng từ $\vec{B}$ vuông góc với trục quay và có độ lớn $\dfrac{\sqrt{2}}{5\pi }$ T. Suất điện động cực đại trong khung dây bằng
	\choice
	{\True $220\sqrt{2}\ V$}
	{220 V}
	{$140\sqrt{2}\ V$}
	{110 V}
	\loigiai{
	\Binhluan{
		\begin{itemize}
			\item
			Tần số góc của khung dây: $\omega =2\pi n=2\pi .50=100\pi $ rad/s.
			\item Suất điện động cảm ứng cực đại: 
			\begin{align*}
				E_0=\omega NBS=100\pi. 500.\dfrac{\sqrt{2}}{5\pi }{.220.10}^{-4}=220\sqrt{2}\ V.
			\end{align*}
		\end{itemize}
	}
	{
	50 vòng/s $\to f=50\ Hz$.	
}
	}
\end{vidu}
\loai{Xác định giá trị hiệu dụng của suất điện động}
\begin{vidu}%[3P3-2.2B][Loai3]
	Một khung dây quay đều quanh trục $\Delta$ trong một từ trường đều có vector cảm ứng từ vuông góc với trục quay. Biết tốc độ quay của khung là 150 vòng/phút. Từ thông cực đại gửi qua khung là $ \dfrac{10}{\pi}\ Wb.$ Suất điện động hiệu dụng trong khung có giá trị là
	\choice
	{25 V}
	{\True $25\sqrt{2}\ V$}
	{50 V}
	{$50\sqrt{2}\ V$}
	\loigiai{
		\Binhluan{
		\begin{itemize}
			\item $n = 150\ \text{vòng/phút}= \dfrac{150}{60}\ \text{vòng/s}= 2,5\ \text{vòng/s} \Rightarrow \omega = 2\pi n = 5\pi$ rad/s.
			\item $E_0=\omega \Phi_0= 50\ V \Rightarrow E=\dfrac{E_0}{\sqrt{2}}=25\sqrt{2} $ V.
	\end{itemize}
}
{Tính suất điện động \textbf{hiệu dụng}}
}
\end{vidu}
\loai{Cho biểu thức từ thông, viết biểu thức suất điện động}
\begin{vidu} %[3P3-1.1B][Loai4]
	Từ thông qua một khung dây có dạng $\Phi =4\cos \left(50\pi t+\dfrac{\pi}{2}\right)\ Wb$. Biểu thức của suất điện động trong khung là
	\choice
	{\True $e=200\pi \sin \left(50\pi t+\dfrac{\pi}{2}\right)\ V$}
	{$e=200\pi \sin \left(50\pi t\right)\ V$}
	{$e=200\pi cos\left(50\pi t+\dfrac{\pi}{2}\right)\ V$}
	{$e=-200\pi \sin \left(50\pi t+\dfrac{\pi}{2}\right)\ V$}
	\loigiai{Suất điện động: $e=-\Phi '=200\pi \sin \left(50\pi t+\dfrac{\pi}{2}\right)\ V$}
\end{vidu}
\loai{Xác định góc tạo bởi vector pháp tuyến và cảm ứng từ tại $t=0$}
\begin{vidu} %[3P3-1.1B][Loai6] 
	\nguon{ĐH - 2011} Một khung dây dẫn phẳng quay đều với tốc độ góc $\omega $ quanh một trục cố định nằm trong mặt phẳng khung dây, trong một từ trường đều có vector cảm ứng từ vuông góc với trục quay của khung. Suất điện động cảm ứng trong khung có biểu thức $e=E_{0}\cos (\omega t+\dfrac{\pi }{2})$ V. Tại thời điểm $t = 0$, vector pháp tuyến của mặt phẳng khung dây hợp với vector cảm ứng từ một góc bằng
	\choice
	{$45^\circ$}
	{\True $180^\circ$}
	{$150^\circ$}
	{$90^\circ$}
	\loigiai{
		\Binhluan{
			\begin{itemize}
				\item Ta có: $e = E_{0}\cos (\omega t+\dfrac{\pi }{2})=E_{0}\sin (\omega t+\pi )$.
				\item So sánh với biểu thức tổng quát: e = $E_{0}\sin (\omega t+\varphi )$, ta có $\varphi =\pi $.
			\end{itemize}
		}
		{Góc $\varphi$ ứng với biểu thức $\Phi$}
	}
\end{vidu}
\loai{Tìm góc tạo bởi mặt phẳng khung và cảm ứng từ} 
\begin{vidu}%[3P3-1.1K][Loai7]
	Một khung dây dẫn kín hình chữ nhật có thể quay đều quanh trục đi qua trung điểm hai cạnh đối diện, trong một từ trường đều có cảm ứng từ $\vec{B},$ vuông góc với trục quay. Suất điện động xoay chiều xuất hiện trong khung có độ lớn cực đại khi mặt phẳng chứa khung dây
	\choice
	{\True song song với $\vec{B}$}
	{vuông góc với $\vec{B}$}
	{tạo với $\vec{B}$ góc $45^\circ$}
	{tạo với $\vec{B}$ góc $60^\circ$}
	\loigiai{
		\Binhluan{
			Ta có: $ \left( \dfrac{e}{E_0} \right)^2+\left( \dfrac{\Phi}{\Phi_0} \right)^2=1 $ $\Rightarrow$ Khi e có độ lớn cực đại $e = E_0$ thì $\Phi = 0$\\ 
			$\Rightarrow$ pha của $\Phi$ là $ \dfrac{\pi}{2}+k\pi $ $\Rightarrow$ $\left( \vec{n},\vec{B} \right)= 90^\circ$.
		}
		{Tìm góc giữa $\vec{B}$ và \textit{mặt phẳng} khung dây}	
	}
\end{vidu}
\loai{Không cho biểu thức từ thông, viết biểu thức suất điện động}
\begin{vidu}%[3P3-1.1B][Loai5]
	Khung dây kim loại phẳng có diện tích $S = 50\ cm^2$ gồm $N = 100$ vòng dây, quay đều với tốc độ 50 vòng/s quanh trục vuông góc với đường sức của một từ trường đều có cảm ứng từ $B = 0,1$ T. Gốc thời gian là lúc vector pháp tuyến của khung dây ngược hướng với vector cảm ứng từ $ \vec{B} $. Từ thông qua khung dây dẫn có biểu thức là
	\choice
	{$\Phi = 0,05\cos (100\pi t)$ Wb}
	{$\Phi = 500\sin(100\pi t + 0,5\pi)$ Wb}
	{\True $\Phi = 0,05\sin(100\pi t - 0,5\pi)$ Wb}
	{$\Phi = 500\cos (100\pi t + 0,5\pi)$ Wb}
	\loigiai{
		\Binhluan{
		\begin{itemize}
			\item Biểu thức từ thông: $\Phi = \Phi_0\cos (\omega t + \varphi )$
			\item $\Phi_0 = NBS = 0,05\ Wb;\ \omega  = 2\pi n = 100\pi$ rad/s.
			\item $t = 0$: $\left| \left( \vec{n},\vec{B} \right) \right|_{t=0}=\left| \varphi
			\right|=\pi \Rightarrow \varphi =\pi$
			$\Rightarrow$ $\Phi = 0,05\cos (100\pi t + \pi ) = 0,05\sin(100\pi t - 0,5\pi)$  Wb.
		\end{itemize}
	}
{$\varphi=\pi$ ứng với biểu thức dạng $\cos$. Ta đổi sang biểu thức dạng $\sin$ cho phù hợp với đáp án}
	}
\end{vidu}
\begin{vidu}%[3P3-1.1K][Loai5]
	Khung dây kim loại phẳng có diện tích $S = 100\ cm^2$ gồm $N = 500$ vòng dây, quay đều với tốc độ 3000 vòng/phút quanh trục vuông góc với đường sức của một từ trường đều có cảm ứng từ $B = 0,1$ T. Gốc thời gian là lúc pháp tuyến của khung dây ngược hướng với vector cảm ứng từ $\vec{B} $. Suất điện động cảm ứng xuất hiện trong khung dây có biểu thức là
	\choice
	{$e = 50\pi \sin(100\pi t - 0,5\pi)$ V}
	{$e = 50\pi \cos (100\pi t - 0,5\pi)$ V}
	{$e = 50\pi \sin(100\pi t + 0,5\pi)$ V}
	{\True $e = 50\pi \cos (100\pi t + 0,5\pi)$ V}
	\loigiai{
		\Binhluan{
		\begin{itemize}	
			\item Biểu thức từ thông: $\Phi = \Phi_0\cos (\omega t + \varphi)$
			\item $\Phi_0 = NBS = 0,5$ Wb; $\omega = 2\pi n = 100\pi$ rad/s.
			\item $\Phi = 0,5\cos (100\pi t + \pi )$ Wb $\Rightarrow$ $e = 50\pi \cos (100\pi t + 0,5\pi )$ V.
		\end{itemize}
}
{Nhớ quy đổi vòng/phút $\to$ vòng/s}
}	
\end{vidu}
\loai{Vận dụng công thức độc lập với thời gian}
\begin{vidu}%[3P3-1.1B][Loai8]
	Một khung dây dẫn quay đều quanh trục xx’ với tốc độ 150 vòng/phút trong một từ trường đều có cảm ứng từ $ \vec{B} $ vuông góc với trục xx’. Ở một thời điểm nào đó thì từ thông gửi qua khung là 4 Wb thì suất điện động cảm ứng trong khung có độ lớn là $15\pi$ V. Từ thông cực đại gửi qua khung là
	\choice
	{\True 5 Wb}
	{$6\pi$ Wb}
	{6 Wb}
	{$5\pi$ Wb}
	\loigiai{
	\Binhluan{
		Ta có: $\omega = 2\pi n = 5\pi$ rad/s $\Rightarrow$ $ \Phi_0=\sqrt{\Phi^2+\dfrac{e^2}{\omega^2}} = 5$ Wb.
	}
{$\Phi$ và $e$ vuông pha nên ta có hệ thức vuông pha}	
}
\end{vidu}
